\section{Mapping TFP Uncertainty to Canonical BCA Wedges}

The mapping implied by the prototype model is summarized in Table
\ref{tab:bca_wedge_mapping}.

\begin{table}[htbp]
\centering
\caption{Canonical BCA Interpretation of TFP Uncertainty}
\label{tab:bca_wedge_mapping}
\begin{tabular}{p{0.34\textwidth}p{0.56\textwidth}}
\toprule
Canonical BCA wedge & Role of TFP uncertainty \\
\midrule
Efficiency wedge \(A_t\) in production
& Direct second-moment shock to \(A_t\) \\
Government/resource wedge \(G_t\) in the resource constraint
& No direct role; absorbs resource residuals and approximation terms \\
Labor wedge \(\tau_{\ell,t}\) in the intratemporal labor FOC
& Indirect, through allocations or misspecification \\
Investment wedge \(\tau_{x,t}\) in the capital Euler equation
& Direct risk correction through \(\Var_t(m^k_{t+1})\) \\
\bottomrule
\end{tabular}
\end{table}

The central lesson is that TFP uncertainty is not a new canonical wedge. It is a
second-moment process attached to the efficiency wedge:
\begin{equation}
    a_{t+1}
    =
    \rho_a a_t
    +
    \exp(V_{A,t})\sigma_A\varepsilon_{A,t+1}.
    \label{eq:efficiency_sv_process}
\end{equation}
However, when agents optimize under this time-varying volatility, the
second-order risk terms can appear as residuals in other equilibrium
conditions. The most natural channel is the investment wedge because the capital
Euler equation contains expectations over future marginal utility and future
capital returns.

A compact accounting representation is
\begin{align}
    a_t
    &=
    \text{efficiency wedge level},
    \label{eq:compact_efficiency}
    \\
    q_t
    &=
    \text{efficiency wedge volatility state},
    \label{eq:compact_efficiency_vol}
    \\
    \xi^x_t
    &\approx
    -\frac{1}{2}\Var_t(m^k_{t+1})
    =
    \lambda_x q_t+\text{higher-order terms},
    \label{eq:compact_investment_loading}
\end{align}
where \(q_t\) is the centered TFP variance state defined in Equation
\eqref{eq:variance_state}. This is the simplest sense in which a second-moment
shock to the efficiency wedge can be represented as an investment wedge in a
first-order BCA accounting exercise.

By contrast, labor and resource wedges should not be mechanically attributed to
TFP uncertainty. They may move in simulations because equilibrium allocations
move, but they are not the first place where conditional variance enters the
prototype model.
